\section*{Informazioni preliminari}
Questo documento presenta il progetto \textit{PetMatch}, un'applicazione web sviluppata dal nostro gruppo per il corso di Tecnologie Web presso l'Università degli Studi di Padova.\\
\textbf{Il sito web è disponibile al link:}
\begin{center}
\url{http://tecweb.studenti.math.unipd.it/acanazza/PetMatch}
\end{center}

\subsubsection*{Credenziali di accesso}
Il sito web consente la visualizzazione come ospite, senza necessità di registrazione. Tuttavia sono presenti \textbf{due aree riservate agli utenti registrati}: l'area utente e l'area amministratore. Per accedere a queste aree, è possibile utilizzare le seguenti credenziali di test:
\begin{center}
\begin{tabular}{|l|c|c|}
\hline
\rowcolor{brown}

\color{white}\textbf{Area} &
\color{white}\textbf{Username} &
\color{white}\textbf{Password} \\
\rowcolor{beige} 
\hline
\textbf{Area utente} & \texttt{user} & \texttt{user} \\
\rowcolor{beige} 
\hline
\textbf{Area amministratore} & \texttt{admin} & \texttt{admin} \\
\hline
\end{tabular}
\end{center}
In entrambi gli account stati predisposti dei dati di esempio per dimostrare le funzionalità del sito.\\
Si nota che queste credenziali sono valide esclusivamente su richiesta esplicita per la consegna del progetto e sono pertanto create appositamente da database; le credenziali create alla registrazione sul sito richiedono di rispettare specifici vincoli di sintassi.
\subsubsection*{Note sulle credenziali}
Si specifica che agli utenti \texttt{user/user} e \texttt{admin/admin} è stata limitata la funzionalità di modifica password e di eliminazione account, per garantire che tutti possano sempre accedere tramite queste credenziali.\\
Solo per scopi di valutazione, è stato creato un ulteriore account di test user, che consente di testare anche le funzionalità di cambio password e di eliminazione.\\
Inoltre, per rendere più chiaro il funzionamento del sito e dimostrare l'esistenza di più account admin, è fornito anche un ulteriore account admin:
\begin{center}
\begin{tabular}{|l|c|c|}
\hline
\rowcolor{brown}
\color{white}\textbf{Area} &
\color{white}\textbf{Username} &
\color{white}\textbf{Password} \\
\rowcolor{beige} 
\hline
\textbf{Area utente} & \texttt{ombretta.gaggi.user@unipd.it} & \texttt{Gaggi!00} \\
\rowcolor{beige} 
\hline
\textbf{Area amministratore} & \texttt{ombretta.gaggi.admin@unipd.it} & \texttt{Gaggi!00} \\
\hline
\end{tabular}
\end{center}
